\section*{Capítulo \ref{ch:b1:lewis-resonance-vsepr} --- Estruturas de Lewis, ressonância e VSEPR}

\begin{solution}{exo:b1:lewis-resonance-vsepr:1}
\ce{H2O}: duas ligações \ce{O-H}, dois \omterm{def:b1:lewis-resonance-vsepr:lewis}{pares não ligantes} em O. \ce{NH3}:
três ligações \ce{N-H}, um \omterm{def:b1:lewis-resonance-vsepr:lewis}{par não ligante} em N. \ce{CO2}: \ce{O=C=O},
dois \omterm{def:b1:lewis-resonance-vsepr:lewis}{pares não ligantes} em cada O. \ce{HCN}: \ce{H-C#N}, um \omterm{def:b1:lewis-resonance-vsepr:lewis}{par não
ligante} em N. \ce{CH2O}: carbono ligado a dois H e com ligação dupla com
O, dois \omterm{def:b1:lewis-resonance-vsepr:lewis}{pares não ligantes} em O. Todas as \omterm{def:b1:lewis-resonance-vsepr:formal-charge}{cargas formais} são nulas nesses
cinco casos. \ce{NH4+}: quatro ligações \ce{N-H}, nenhum \omterm{def:b1:lewis-resonance-vsepr:lewis}{par não ligante};
nitrogênio $5 - 0 - 4 = +1$, a carga do íon.
\end{solution}

\begin{solution}{exo:b1:lewis-resonance-vsepr:2}
\ce{H2S}: \ce{AX2E2}, angular, abaixo de $109.5^\circ$ (medido: $92^\circ$).
\ce{PCl3}: \ce{AX3E}, pirâmide trigonal, cerca de $100^\circ$.
\ce{BF3}: \ce{AX3}, trigonal plana, $120^\circ$. \ce{SiH4}: \ce{AX4},
tetraédrica, $109.5^\circ$. \ce{CO3^{2-}}: \ce{AX3} (a ligação dupla
conta uma vez), trigonal plana, $120^\circ$.
% ledger: geo:H2S.aHSH, geo:PCl3.aClPCl
\end{solution}

\begin{solution}{exo:b1:lewis-resonance-vsepr:3}
Apolares: \ce{CCl4} (\ce{AX4} tetraédrica), \ce{BF3} (\ce{AX3}),
\ce{CO2} (\ce{AX2} linear): os \omterm{def:b1:lewis-resonance-vsepr:dipole}{dipolos de ligação} se anulam. Polares:
\ce{CH2Cl2} (dois tipos diferentes de ligação em um tetraedro), \ce{NH3}
(piramidal), \ce{SO2} (angular).
\end{solution}

\begin{solution}{exo:b1:lewis-resonance-vsepr:4}
\ce{H+} é o \omterm{def:b1:lewis-resonance-vsepr:lewis-acid}{ácido de Lewis} e \ce{H2O} a base: a nova ligação \ce{O-H} de
\ce{H3O+} é dativa no momento em que se forma (depois, as três ligações
\ce{O-H} são idênticas). \ce{Cu^{2+}} é o ácido e cada \ce{NH3} uma base:
as quatro ligações \ce{Cu-N} são dativas.
\end{solution}

\begin{solution}{exo:b1:lewis-resonance-vsepr:5}
\chemfig{CH_3-C(=[:60]O)-[:-60]\chemabove{O}{\scriptstyle\ominus}}
$\leftrightarrow$
\chemfig{CH_3-C(-[:60]\chemabove{O}{\scriptstyle\ominus})=[:-60]O}.
As duas estruturas são equivalentes: as duas ligações \ce{C-O} têm ordem
1.5 e o mesmo comprimento, entre o da ligação dupla e o da simples. No
ácido metanoico, os dois oxigênios diferem (um leva o H), as estruturas
não são equivalentes e as ligações mantêm comprimentos distintos, 120.2
(\ce{C=O}) e
\qty{134.3}{pm} (\ce{C-OH}).
\end{solution}

\begin{solution}{exo:b1:lewis-resonance-vsepr:6}
$N = 6 + 4 \times 6 + 2 = 32$ elétrons. Com quatro ligações simples e três
\omterm{def:b1:lewis-resonance-vsepr:lewis}{pares não ligantes} em cada oxigênio: enxofre $6 - 0 - 4 = +2$, cada
oxigênio $6 - 6 - 1 = -1$; total $+2 - 4 = -2$. Forma \ce{AX4},
tetraédrica. Com duas ligações \ce{S=O}, a escolha de quais dois dos
quatro oxigênios fazem ligação dupla dá $\binom{4}{2} = 6$ estruturas.
\end{solution}

\begin{solution}{exo:b1:lewis-resonance-vsepr:7}
\ce{SF4}: \ce{AX4E}, gangorra. \ce{ClF3}: \ce{AX3E2}, em T.
\ce{XeF2}: \ce{AX2E3}, linear. \ce{XeF4}: \ce{AX4E2}, quadrado plana. Em
uma bipirâmide trigonal, uma posição axial tem três vizinhos a
$90^\circ$, uma posição equatorial apenas dois: os volumosos \omterm{def:b1:lewis-resonance-vsepr:lewis}{pares não
ligantes} vão para onde encontram menos repulsões a $90^\circ$.
\end{solution}

\begin{solution}{exo:b1:lewis-resonance-vsepr:8}
O enxofre e o fósforo são maiores e menos eletronegativos que o oxigênio
e o nitrogênio: os \omterm{def:b1:lewis-resonance-vsepr:lewis}{pares ligantes} ficam mais longe do átomo central e são
puxados em direção ao hidrogênio, de modo que se repelem menos, e os \omterm{def:b1:lewis-resonance-vsepr:lewis}{pares
não ligantes} comprimem as ligações em direção a $90^\circ$.
\end{solution}

\begin{solution}{exo:b1:lewis-resonance-vsepr:9}
\ce{CH3-C#N}: o carbono do \ce{CH3} é sp$^3$, o carbono da nitrila sp, o
nitrogênio sp; 5 ligações $\sigma$ (três \ce{C-H}, \ce{C-C}, uma em
\ce{C#N}) e 2 ligações $\pi$. \ce{CH2=CH-CH3}: os dois carbonos do alceno
são sp$^2$, o carbono metílico sp$^3$; 8 ligações $\sigma$ (seis \ce{C-H},
duas \ce{C-C}) e 1 ligação $\pi$.
\end{solution}

\begin{solution}{exo:b1:lewis-resonance-vsepr:10}
As duas moléculas são piramidais, com um \omterm{def:b1:lewis-resonance-vsepr:lewis}{par não ligante} no nitrogênio.
O \omterm{def:b1:lewis-resonance-vsepr:lewis}{par não ligante} dá uma contribuição orientada (do negativo para o
positivo) do par em direção ao núcleo. Em \ce{NH3}, o nitrogênio é a
extremidade negativa de cada ligação: os \omterm{def:b1:lewis-resonance-vsepr:dipole}{dipolos de ligação} apontam de N
para os hidrogênios, no mesmo sentido da contribuição do \omterm{def:b1:lewis-resonance-vsepr:lewis}{par não ligante},
e se somam. Em \ce{NF3}, o flúor é a extremidade negativa: os \omterm{def:b1:lewis-resonance-vsepr:dipole}{dipolos de
ligação} apontam dos flúores para N, contra a contribuição do \omterm{def:b1:lewis-resonance-vsepr:lewis}{par não
ligante}, e os dois quase se anulam.
\end{solution}

\begin{solution}{exo:b1:lewis-resonance-vsepr:11}
$\mu_b = \mu/(2\cos(\theta/2)) = 1.63/(2\cos 59.75^\circ) =
\qty{1.62}{\debye} = \qty{5.40e-30}{C.m}$. Então $\delta = \mu_b/(e\,d) =
5.40 \times 10^{-30}/(1.602 \times 10^{-19} \times 143.2 \times 10^{-12})
= 0.24$.
\end{solution}

\begin{solution}{exo:b1:lewis-resonance-vsepr:12}
\ce{N=N=O} com \omterm{def:b1:lewis-resonance-vsepr:formal-charge}{cargas formais} $-1$, $+1$, $0$; \ce{N#N-O} com $0$,
$+1$, $-1$. A ligação \ce{N-N} (\qty{112.8}{pm}) é próxima da ligação
tripla de \ce{N2} (\qty{109.8}{pm}): a estrutura \ce{N#N-O} pesa mais, e
ela põe a carga negativa no oxigênio, o átomo mais eletronegativo. A
ligação \ce{N-O}, porém, é mais curta que uma ligação simples: a outra
estrutura também contribui.
\end{solution}

\begin{solution}{pb:b1:lewis-resonance-vsepr:1}
\textbf{1.} \ce{N2} 10, \ce{NO} 11, \ce{NO2} 17, \ce{N2O} 16, \ce{O3} 18,
\ce{HNO3} 24.
\textbf{2.} \ce{N#N} com um \omterm{def:b1:lewis-resonance-vsepr:lewis}{par não ligante} em cada nitrogênio.
\textbf{3.} \ce{N=O} com dois \omterm{def:b1:lewis-resonance-vsepr:lewis}{pares não ligantes} no oxigênio, um \omterm{def:b1:lewis-resonance-vsepr:lewis}{par não
ligante} e um elétron isolado no nitrogênio. Com 11 elétrons, número
ímpar, um átomo sempre fica com uma contagem ímpar: o nitrogênio tem 7
elétrons em torno dele.
\textbf{4.} \ce{N=N=O}: $-1$, $+1$, $0$; \ce{N#N-O}: $0$, $+1$, $-1$.
\textbf{5.} \ce{O=O-O}: oxigênio central $+1$ (um \omterm{def:b1:lewis-resonance-vsepr:lewis}{par não ligante}, três
ligações), o oxigênio terminal com ligação simples $-1$, o outro 0.
\textbf{6.} Como no \cref{ex:b1:lewis-resonance-vsepr:hno3}: nitrogênio $+1$,
o oxigênio com ligação simples e sem hidrogênio $-1$, os outros 0.
\textbf{7.} \ce{NO} e \ce{NO2} (11 e 17 elétrons).
\textbf{8.} Duas estruturas equivalentes, com a ligação dupla de um lado
ou do outro: ordem 1.5 para cada ligação.
\textbf{9.} Sim: \qty{127.8}{pm} fica entre a ligação dupla
(\qty{120.8}{pm}) e a ligação simples (\qty{147.5}{pm}).
\textbf{10.} A ligação dupla em cada um dos três oxigênios, por vez; cada
ligação \ce{N-O} é dupla em uma estrutura em três: ordem $4/3$.
\textbf{11.} A ligação longa, \qty{140.6}{pm}, é \ce{N-OH}, uma ligação
simples. Os dois oxigênios sem hidrogênio compartilham uma ligação $\pi$
por meio de duas \omterm{def:b1:lewis-resonance-vsepr:resonance}{estruturas de ressonância} equivalentes: duas ligações de
ordem 1.5, 121.1 e
\qty{119.9}{pm}.
\textbf{12.} O nitrogênio leva um elétron isolado em vez de um \omterm{def:b1:lewis-resonance-vsepr:lewis}{par não
ligante}: ele repele os \omterm{def:b1:lewis-resonance-vsepr:lewis}{pares ligantes} menos do que um par repeliria, e
as ligações se abrem além de $120^\circ$.
\textbf{13.} \ce{N#N-O}, com a carga negativa no oxigênio.
\textbf{14.} \ce{O3}: \ce{AX2E}, angular. \ce{N2O}: \ce{AX2} (N central),
linear. \ce{NO3-}: \ce{AX3}, trigonal plana. Nitrogênio de \ce{HNO3}:
\ce{AX3}, trigonal plana.
\textbf{15.} O \omterm{def:b1:lewis-resonance-vsepr:lewis}{par não ligante} do oxigênio central repele os \omterm{def:b1:lewis-resonance-vsepr:lewis}{pares
ligantes} mais do que eles se repelem entre si.
\textbf{16.} O ozônio é angular e seu átomo central difere dos átomos
das extremidades (\omterm{def:b1:lewis-resonance-vsepr:formal-charge}{carga formal} positiva no meio, carga negativa
compartilhada pelas extremidades): um pequeno dipolo ao longo da
bissetriz. \ce{N2O} é linear, mas suas duas extremidades são átomos
diferentes: um pequeno dipolo. \ce{N2} é simétrico: nenhum.
\textbf{17.} \ce{CO2} é linear (\ce{AX2}): seus \omterm{def:b1:lewis-resonance-vsepr:dipole}{dipolos de ligação} se
anulam. \ce{SO2} é angular (\ce{AX2E}): eles se somam.
\textbf{18.} \ce{AX3E}: o \omterm{def:b1:lewis-resonance-vsepr:lewis}{par não ligante} fecha os ângulos abaixo de
$109.5^\circ$; medido: $106.7^\circ$.
\textbf{19.} \ce{AX2E2}: dois \omterm{def:b1:lewis-resonance-vsepr:lewis}{pares não ligantes} o fecham ainda mais;
$104.5^\circ$.
\textbf{20.} $\mu = 2\mu_{OH}\cos(\theta/2)$
(\cref{prop:b1:lewis-resonance-vsepr:dipole-sum}).
\textbf{21.} $\mu_{OH} = 1.857/(2\cos 52.24^\circ) = 1.857/1.2246 =
\qty{1.52}{\debye}$.
\textbf{22.} $1.516 \times 3.336 \times 10^{-30} = \qty{5.06e-30}{C.m}$.
\textbf{23.} $e\,d = 1.602 \times 10^{-19} \times 95.8 \times 10^{-12} =
\qty{1.535e-29}{C.m}$, ou seja, \qty{4.60}{\debye}.
\textbf{24.} $\delta = 5.06 \times 10^{-30}/1.535 \times 10^{-29} =
\textbf{0.33}$: a ligação \ce{O-H} da água porta um terço da carga que
uma ligação totalmente iônica portaria.
\end{solution}
